\documentclass[11pt,a4paper]{article}
\usepackage[T1]{fontenc}
\usepackage[utf8]{inputenc}
\usepackage{lmodern}
\usepackage[margin=2.4cm]{geometry}
\usepackage{booktabs}
\usepackage{siunitx}
\usepackage{xcolor}
\usepackage{listings}
\usepackage[hidelinks]{hyperref}
\usepackage{microtype}

\sisetup{detect-all}
\lstset{basicstyle=\ttfamily\small, columns=fullflexible, breaklines=true, frame=none, xleftmargin=1em}
\newcommand{\code}[1]{\texttt{#1}}

\title{Two decode kernels for gfx906 in llama.cpp:\\a quantized matvec and flash decoding on the Radeon VII}
\author{Hermes Reisner}
\date{November 2026}

\begin{document}
\maketitle

\begin{abstract}
\noindent
Generating one token with a quantized language model is a chain of matrix-vector products and, in long chats, an
attention pass over the whole cache. On AMD's gfx906 (Radeon VII, Instinct MI50 and MI60) llama.cpp's generic HIP
kernels use the card's memory bandwidth poorly for both. This note describes two kernels for single-token decode: a
matrix-vector product for Q4\_0, Q4\_1 and IQ4\_NL weights, and a flash-decoding attention kernel for a q8\_0 cache. On a Radeon VII the matvec is 2 to 34\,\% faster than upstream on dense products with at least 512 rows
(median 14 to 18\,\% depending on the format) and the attention kernel needs 453 instead of 787\,\textmu s at 64k keys.
In whole models on the Radeon VII alone, Gemma 4 26B (Q4\_0) writes 10\,\% faster and Qwen 3.8 27B 3\,\% faster;
prompt processing and perplexity do not change.
The code was written with an AI coding agent (Claude Code). Every number here was measured on the card, every kernel
is checked against llama.cpp's CPU reference, and the cases where the kernels did not help are reported too. Similar
ideas exist in community forks; Section~\ref{sec:related} says which.
\end{abstract}

\section{Setup and motivation}

The machine is a home server with a Radeon VII (gfx906, 60 compute units, 16\,GB HBM2 at 1\,TB/s, PCIe 3.0 x4) and an
RX 9060 XT (gfx1200, 16\,GB) on Fedora 44 with Fedora's ROCm 7.1 packages. It serves a chat site, and the model this
work was done for is Inkling 2 Mini (Thinking Machines): a 26.4B mixture of experts with 128 experts of which 8 are used
per token, hidden size 2304, expert FFN size 1152, 24 query heads on 8 key/value heads of size 128, and 24 stored layers
of which layers 6 to 17 run three times, so a token passes through 48 layers. Stored layers 0 to 17 live on the Radeon
VII and the rest on the RX 9060 XT.

At batch size one, decode is a long chain of small kernels. A profile showed about 2,200 kernel launches per token, and
the VII's matrix-vector products reached only 200 to 260 of its roughly 1000\,GB/s. Raising power limits changed
nothing, raising the memory clock changed nothing, and raising the core clock helped a little, so the kernels were
bound by latency and instruction issue rather than by bandwidth. That is where both kernels start.

\paragraph{gfx906 in brief.} A compute unit has four 16-wide SIMDs that run 64-wide wavefronts, up to 10 per SIMD
(40 per compute unit), and 64\,KB of LDS. The chip has integer dot-product instructions
(\code{v\_dot4\_i32\_i8}), which llama.cpp uses as \code{dp4a}. The MI50 is the same chip with 16 or 32\,GB, the MI60
has 64 compute units and 32\,GB. Only the Radeon VII was tested.

\section{The quantized matvec}

\subsection{What upstream does}

llama.cpp quantizes the activation vector to q8\_1 (36-byte blocks of 32 values: a scale and a sum as \code{half2},
then 32 \code{int8}) and runs one kernel family, \code{mul\_mat\_vec\_q}, for all weight formats. Each thread handles
\code{vdr} 32-bit words of a block, so for the four formats here two threads share one 32-value block; on GCN a row
gets one 128-thread workgroup. The weight blocks are 18 bytes (Q4\_0, IQ4\_NL), 20 bytes (Q4\_1) and 34 bytes (Q8\_0),
so most of them are only 2-byte aligned.

\subsection{Whole blocks per lane}

In the new kernel each lane computes whole blocks and accumulates them in a register; lanes of a row only meet in a
final shuffle reduction. A Q4\_0 or IQ4\_NL block is read with nine 2-byte loads (the block is only 2-byte aligned)
and a Q4\_1 block, which is 4-byte aligned, with five 4-byte loads. An earlier version also covered Q8\_0 (34-byte
blocks, seventeen 2-byte loads); it was not faster than upstream and was dropped. Eight \code{dp4a}
instructions then cover the 32 values. IQ4\_NL needs a 16-entry table lookup per nibble; it uses the same
\code{v\_perm}-based helper as upstream (\code{get\_int\_from\_table\_16}) and is still the most expensive format on
this card: memory at 1200 instead of 1000\,MHz changed nothing, the core clock did. This is why the Inkling model file
uses Q4\_1 for the experts.

The microbenchmark in \code{bench/mvbench.hip} compares this ``direct'' variant with others on the Inkling expert
shapes: staging the weights in LDS with 16-byte loads first, several blocks' loads in flight per lane, and aligned
16-byte plus 4-byte loads with \code{alignbyte}. All were slower than the direct version.

\subsection{Experts}

For \code{MUL\_MAT\_ID} (mixture-of-experts products) a 128-thread workgroup computes 16 rows with 8 lanes each and
first copies the activation column into LDS, so the 16 rows share one copy instead of reading it from memory 16 times.
The activation must fit (at most 1024 blocks, 36\,KB). Q8\_0 experts are left to the generic kernel, which was faster
there (103 against 121\,\textmu s).

\subsection{Dense products and a floor at 2048 workgroups}
\label{sec:floor}

The first version used 64-thread workgroups (one wavefront) with 2 rows of 32 lanes. It was faster on large products
and on Inkling's shapes, but on Qwen 3.8 27B in Q4\_0 the whole model got 5\,\% slower. A sweep over the decode shapes
of several models and eight lane layouts explained it. Every 64-thread layout has a time floor of about 57\,\textmu s when the grid has about 2048
workgroups, whatever the matrix size: 2048 rows at 1 row per workgroup, 4096 rows at 2, 8192 at 4 and 16384 at 8 all
land there (Table~\ref{tab:floor}). Grids slightly above (2304 workgroups) or below (1536) are fine, and a matrix large
enough to need more than 57\,\textmu s anyway hides it. Upstream's 128-thread workgroups do not show it, and neither
do the same layouts with two wavefronts per workgroup. The cause is not known; it looks like a limit in how the card
dispatches many one-wavefront workgroups.

\begin{table}[ht]
\centering
\caption{Q4\_0 matvec, one token, Radeon VII, \textmu s. Layouts are rows per workgroup x lanes per row.}
\label{tab:floor}
\begin{tabular}{rrr rrrr rr}
\toprule
& & & \multicolumn{4}{c}{64 threads} & \multicolumn{2}{c}{128 threads} \\
\cmidrule(lr){4-7}\cmidrule(lr){8-9}
rows & K & upstream & 1x64 & 2x32 & 4x16 & 8x8 & 2x64 & 4x32 \\
\midrule
 2048 & 4096 & 14.9 & \textbf{57.6} & 13.6 & 16.3 & 19.9 & 13.6 & 13.9 \\
 4096 & 2048 & 15.9 & 15.3 & \textbf{57.6} & 13.6 & 16.8 & 14.4 & 13.5 \\
 4096 & 4096 & 23.6 & 24.8 & \textbf{57.5} & 21.2 & 28.6 & 20.6 & 21.9 \\
 4608 & 2048 & 16.9 & 16.3 & 14.0 & 14.7 & 17.7 & 15.6 & 14.3 \\
 8192 & 2048 & 25.2 & 24.2 & 24.6 & \textbf{57.6} & 21.8 & 22.5 & 20.5 \\
16384 & 2048 & 45.3 & 42.8 & 36.9 & 48.9 & \textbf{57.5} & 38.8 & 34.4 \\
\bottomrule
\end{tabular}
\end{table}

The final kernel therefore uses 128-thread workgroups everywhere: 4 rows of 32 lanes, 2 rows of 64 lanes below 2048
rows, 8 rows of 16 lanes when $K < 2048$, and the generic kernel below 512 rows, where it is as fast. The rule was
chosen on the same shapes it is reported on below, so those numbers are not an independent test; the whole-model
runs are. Dense Q8\_0 was not faster with any layout (median $+2\,\%$), so Q8\_0 stays with the generic kernel.

\subsection{Fused gate}

Upstream fuses the gate and up products of a feed-forward layer with the activation function (SwiGLU, GeGLU and their
clamped variants) into one matvec. The first version skipped fused products, so on Gemma 4 and Qwen it only reached
the attention and down projections, and Gemma 4 26B gained less than 1\,\%. The kernel now computes the gate row next
to the up row and applies the activation with upstream's \code{ggml\_cuda\_op\_glu\_single}. Products with a fused
bias or scale still go to the generic kernel.

\subsection{Results}

Table~\ref{tab:mmv} shows selected shapes; over all 49 shapes with at least 512 rows the change against upstream is
$-2$ to $-28\,\%$ for Q4\_0 (median $-14\,\%$), $-5$ to $-34\,\%$ for Q4\_1 (median $-18\,\%$) and 0 to $-31\,\%$ for
IQ4\_NL (median $-15\,\%$). Below 512 rows the generic kernel runs, so nothing changes there.

\begin{table}[ht]
\centering
\caption{Matrix-vector product for one token on the Radeon VII, \textmu s, upstream / this kernel
(\code{test-backend-ops perf}, llama.cpp \code{94f3af3d0}). Rows x K; the example names a model with that shape.}
\label{tab:mmv}
\small
\begin{tabular}{rl rr rr rr}
\toprule
shape & example & \multicolumn{2}{c}{Q4\_0} & \multicolumn{2}{c}{Q4\_1} & \multicolumn{2}{c}{IQ4\_NL} \\
\midrule
1024 x 4096 & Llama 3 8B, k/v & 9.9 & 9.1 & 9.3 & 8.2 & 10.4 & 9.4 \\
2304 x 2304 & Inkling 2 Mini & 12.5 & 11.7 & 12.4 & 9.6 & 14.0 & 12.8 \\
4096 x 4096 & Llama 3 8B, q/o & 23.6 & 21.9 & 21.8 & 20.5 & 24.4 & 22.9 \\
4096 x 14336 & Llama 3 8B, down & 71.6 & 64.2 & 65.4 & 55.2 & 73.6 & 67.0 \\
14336 x 4096 & Llama 3 8B, up & 68.2 & 56.1 & 61.6 & 53.1 & 70.5 & 57.8 \\
6144 x 5120 & Qwen 3.8 27B & 39.6 & 34.5 & 36.5 & 30.6 & 40.7 & 34.7 \\
5120 x 17408 & Qwen 3.8 27B, down & 106.4 & 88.4 & 97.4 & 76.3 & 109.7 & 90.9 \\
8192 x 8192 & 70B, q/o & 78.2 & 66.2 & 71.2 & 64.0 & 80.5 & 68.3 \\
28672 x 8192 & 70B, up & 259.0 & 211.3 & 234.2 & 200.3 & 266.3 & 216.8 \\
151936 x 2048 & output head, 2048 hidden & 378.7 & 272.9 & 377.4 & 247.9 & 410.6 & 285.3 \\
\bottomrule
\end{tabular}
\end{table}

In whole models (Table~\ref{tab:models}) the gain is smaller, because other operations take their share of each
token: the Q6\_K output head of Qwen 3.8, its recurrent layers and the attention. The Qwen file in IQ4\_XS, which the
kernel does not cover, is the control and does not change.

\begin{table}[ht]
\centering
\caption{Writing speed (\code{llama-bench}, 128 tokens, flash attention on), Radeon VII alone, tok/s, mean of the runs
after the first. The Qwen files in Q4\_0 and IQ4\_NL were requantized from the Q6\_K file for this test.}
\label{tab:models}
\begin{tabular}{lrrr}
\toprule
model & upstream & this kernel & change \\
\midrule
Gemma 4 26B-A4B, QAT Q4\_0 & 91.9 & 101.5 & $+10\,\%$ \\
Qwen 3.8 27B, Q4\_0 & 30.1 & 31.1 & $+3\,\%$ \\
Qwen 3.8 27B, IQ4\_NL & 28.5 & 29.3 & $+3\,\%$ \\
Qwen 3.8 27B, UD-IQ4\_XS (not covered) & 25.4 & 25.6 & $\pm 0$ \\
\bottomrule
\end{tabular}
\end{table}

Output and prompt speed were checked too. Perplexity on wikitext-2 with batches of one token, so that every token goes
through the kernel, is 6.031 with the generic kernel and 6.034 with this one (Qwen 3.8 27B Q4\_0, 4 chunks of 512).
Greedy answers start identically and drift apart after 10 to 40 tokens, which is what a different summation order
does at near-ties; both versions stay coherent. Prompt processing (512 tokens) is 350.7 against 349.2\,tok/s, i.e.
unchanged, since the kernel only runs for one token.

\section{Flash decoding}

\subsection{The vector kernel at long context}

For one query token llama.cpp uses its vector flash-attention kernel. It assigns one query head per workgroup and can
split the keys over several workgroups (\code{parallel\_blocks}) with a combine pass. With 3 query heads per key/value
head (Inkling 2: 24 on 8), each key/value row is therefore read three times, and on the VII the kernel moved about
160\,GB/s at 32k to 64k keys with a q8\_0 cache.

\subsection{Design}

The new kernel (\code{flash-decoding/fattn-gcn-dec.cu}) is specialised to head size 128, 3 query heads per key/value
head and a q8\_0 cache. The grid has one workgroup per key/value head and chunk of keys; a workgroup handles all 3 query
heads, so it reads each key and value row once.
\begin{enumerate}
\item The 3 query vectors are quantized to \code{int8} per 32 values into LDS, with the softmax scale folded in.
\item The workgroup walks over its chunk in tiles of 256 keys. Each lane loads one key's q8\_0 row (136 bytes, with
16-byte loads) and computes its 3 scores with \code{dp4a}, then the optional Inkling 2 relative-position bias is added.
The next tile's rows are loaded while the current tile does steps 3 and 4.
\item An online softmax keeps a running maximum and sum per head; the accumulators are rescaled when the maximum grows.
\item For the values, 8 lanes cover one row (16 values each), 32 rows are processed at a time and 4 rows of loads are
kept in flight. Accumulation is in \code{half2}.
\item Each chunk writes its partial output and its (maximum, sum); a second small kernel combines the chunks with a
log-sum-exp and adds the attention sinks.
\end{enumerate}
The number of tiles per workgroup was measured: about $\sqrt{t}/2$ for $t$ tiles of keys (2 at 4k keys, 4 at 16k,
8 at 64k). Fewer leave compute units idle, more lengthen each workgroup's chain of memory round trips.

Two things did not work first. A version with \code{\_\_launch\_bounds\_\_(256, 2)} was limited to 128 registers and
spilled 112 bytes per lane to scratch, which made it slower than the vector kernel; with \code{(256, 1)} it uses 155 to
202 registers and does not spill (checked with \code{-Rpass-analysis=kernel-resource-usage}). And a version without the
tile loop, one tile per workgroup, was slow at long context because it launched too many short workgroups.

\subsection{Results}

\begin{table}[ht]
\centering
\caption{Attention for one token, 24 query heads on 8 key/value heads, q8\_0 cache, Inkling 2 bias, Radeon VII
(\code{test-backend-ops perf}).}
\label{tab:fa}
\begin{tabular}{rrrr}
\toprule
keys & vector kernel (\textmu s) & this kernel (\textmu s) & change \\
\midrule
 1k  & 23.7 & --- & \\
 4k  & 64.1 & --- & \\
 8k  & 110.2 & 92.4 & $-16\,\%$ \\
16k  & 210.4 & 162.0 & $-23\,\%$ \\
32k  & 407.5 & 225.7 & $-45\,\%$ \\
64k  & 786.7 & 453.1 & $-42\,\%$ \\
\bottomrule
\end{tabular}
\end{table}

Below 6144 keys the vector kernel is as fast, so the new kernel is only used from there on
(\code{GGML\_CUDA\_GCN\_FA\_DEC\_MIN\_KV}). In Inkling 2 most layers use a 1024-token sliding window and never reach
that, so the kernel only changes the few global layers. In the whole model, writing 160 tokens after a document of 34k
tokens went from 38.5 to 40.4\,tok/s and after 68.5k tokens from 33.5 to 37.0 (mean of two runs, 131k context, both
cards otherwise idle). After 8.6k tokens the difference (43.9 against 43.5) is within the run-to-run noise. Reading
speed does not change.

\section{How it was checked}

\begin{itemize}
\item \textbf{Against the CPU.} \code{test-backend-ops} compares each GPU op with llama.cpp's CPU implementation.
Upstream llama.cpp with the matvec patch, on the Radeon VII: 683 \code{MUL\_MAT} and \code{MUL\_MAT\_ID} cases and 341
fused gate/up cases (\code{MUL\_MAT\_FUSION}) for Q4\_0, Q4\_1, IQ4\_NL and Q8\_0. The Inkling 2 tree, with both
cards in one process: on the Radeon VII 431 \code{MUL\_MAT} and 262 \code{MUL\_MAT\_ID} cases, 67 Inkling 2 op cases
including attention with the bias up to 33k keys, and 6 added plain attention cases (7000 to 16384 keys, with and
without sinks); on the RX 9060 XT 914 float \code{MUL\_MAT} cases without rocBLAS and the 67 Inkling 2 cases. All pass.
\item \textbf{The whole model.} Inkling 2 Mini finds the needle in a 64k and a 118k token document; wikitext-2
perplexity of the model file is 6.982 against 6.916 for ggml-org's IQ4\_NL file; 12 of 12 tasks of a small private
evaluation pass.
\item \textbf{Speed.} Every speed number is from \code{test-backend-ops perf}, \code{llama-bench} or the chat server on
the Radeon VII at stock clocks (1801\,MHz, memory 1000\,MHz), with the other card idle unless stated.
\end{itemize}
Not checked: an MI50 or MI60 (same instruction set, the MI60 has 64 compute units), other ROCm versions than Fedora's
7.1, other operating systems, and models with other head sizes or group sizes for the attention kernel.

\section{Related work}
\label{sec:related}

Both ideas have precedents, and the gfx906 community has been active:
\begin{itemize}
\item \textbf{iacopPBK/llama.cpp-gfx906} already computes whole Q4\_0, Q4\_1 and Q8\_0 blocks per lane at one column,
with 2 rows of 32 lanes, but only for $K \le 1024$; its comment says it wins only on small matrices. It has no IQ4\_NL
path in that kernel.
\item \textbf{exabit-io/mx-llama.cpp} (and the retired exabit-io/llama.cpp) has a gfx906 Q8\_0 path that processes a
whole block per thread at one column, with an optional LDS-staged activation.
\item \textbf{Martin1701/MI50-llama.cpp} has a gfx906 attention kernel that shares a q8\_0 key/value load across six
query heads, quantizes the query into shared memory and uses upstream's split and combine, specialised to head size
256 (Qwen 3.8 27B). It accumulates in \code{float}. \textbf{0c33/mi50-vllm-lab} has a vLLM kernel with 256-key tiles,
six-head sharing and an online softmax for vLLM's own int8 cache.
\item \textbf{Upstream} has merged GCN tuning for the prompt-processing kernels (\#27841, \#30021, \#30022) and the
tile attention kernel (\#15769, \#15884, \#15927); \#29534 (wave64 for GCN q8 attention) and \#30344 were open in
November 2026. Issue \#31765 proposed gfx906 decode kernels and is closed.
\end{itemize}
What this note adds: the matvec for Q4\_0, Q4\_1 and IQ4\_NL at every $K$, with the expert path, the fused gate and a
layout rule measured on 51 shapes per format; the observation about 2048 one-wavefront workgroups (Section~\ref{sec:floor}); an attention
kernel for head size 128 with 3 query heads per key/value head; and a patch against current upstream llama.cpp rather
than a fork.

\section{What did not help}

\begin{itemize}
\item Speculative decoding with Inkling's DSpark draft model: at most +4\,\% on short answers and slower on long ones,
because verifying 4 tokens touches up to 32 experts.
\item Raising power limits, overclocking the memory, other output-head formats, other layer splits.
\item Staging weights in LDS, aligned 4-byte loads for the 18-byte blocks, a dedicated expert-router kernel (the fused
one takes 1.6\,\textmu s), polling instead of interrupts (\code{HSA\_ENABLE\_INTERRUPT=0}, much slower).
\item A GQA-packed version of the generic vector attention kernel and 64-thread blocks for it: correct, but slower
below 6k keys.
\item Fusing Inkling's conv state store with its short convolution: correct, no measurable gain.
\item A Vulkan fork that reaches 108\,tok/s on one RX 7900 XT: 32\,tok/s here, the Radeon VII is slow under RADV.
\end{itemize}

\section{Limits}

The matvec covers one token and the formats Q4\_0, Q4\_1 and IQ4\_NL, not Q8\_0 and not the K-quants that most
downloaded files use. The attention kernel only covers head size 128, 3 query heads per key/value head and a q8\_0 cache, and it
lives in the Inkling 2 tree, since upstream llama.cpp does not support Inkling 2 yet. Prompt processing uses the stock
kernels. The matvec patch is not upstreamed; llama.cpp asks contributors to disclose AI use and to be able to explain
and maintain every line, which is a decision for the author, not for the agent that wrote it.

\appendix
\section{Reproducing}

\begin{lstlisting}
git clone https://github.com/ggml-org/llama.cpp && cd llama.cpp
git checkout 94f3af3d0            # the commit the patch was tested on
git apply ../matvec/gcn-mmv.diff
HIPCXX="$(hipconfig -l)/clang" HIP_PATH="$(hipconfig -R)" cmake -S . -B build \
    -DGGML_HIP=ON -DAMDGPU_TARGETS=gfx906 -DCMAKE_BUILD_TYPE=Release
cmake --build build -j 8 --target llama-bench test-backend-ops
build/bin/test-backend-ops test -o MUL_MAT,MUL_MAT_ID -p "type_a=(q4_0|q4_1|iq4_nl|q8_0)"
GGML_CUDA_GCN_MMV=0 build/bin/llama-bench -m model.gguf -p 0 -n 128   # upstream kernel
GGML_CUDA_GCN_MMV=1 build/bin/llama-bench -m model.gguf -p 0 -n 128   # this one (default)
\end{lstlisting}
For the Inkling 2 tree with the attention kernel, see \code{inkling2/README.md} in the repository.

\end{document}
